\section{Introduction}\label{sec:intro}

On \dv{lean.session.date_long} we rented a cloud server (a pod) with two NVIDIA RTX PRO 6000 GPUs, 96\,GB each, and
asked it to load a version of DeepSeek-V4.1-Flash whose routed experts are quantized to \dv{cfg.exl3_bits} bits, with
every expert in GPU memory. The load failed with a CUDA out-of-memory error, and so did the offload build with an
empty plan. With a plan that moves \dv{static.gen30.host_target_share_pct}\% of the routed experts to the host's
ordinary memory, a back room that the GPUs can still reach, the same offload build started answering requests.

DeepSeek publishes the model in \dv{static.orig.total_gib}\,GiB, and the community pack with 3-bit experts takes
\dv{static.pack.total_gib}\,GiB. Not all of it has to be on the GPUs: the two Engram lookup tables and the token
embedding (\dv{static.noplan.host_gib}\,GiB together) are kept in the host's memory before any expert moves,
and a vision encoder is not loaded at all. The rest, \dv{static.noplan.vram_gib}\,GiB, must be split between the two cards, which comes to
\dv{static.noplan.vram_per_gpu_gib}\,GiB per card, and the driver reports \dv{expl.gpu_total_gib}\,GiB per card. The weights alone are \dv{expl.noplan_over_per_gpu_gib}\,GiB
too large for each GPU, before any memory is set aside for the KV cache, which holds the context of every conversation
in progress.

The software that turned the failed boot into a working one was written by other people. Section~\ref{sec:attribution}
names them and what each one built. This paper is about what that software buys and what it costs, and the question we
set ourselves was:
\begin{quote}
\emph{What exactly does host-memory offload make possible on two 96\,GB GPUs, and at what cost in speed, capacity and
quality?}
\end{quote}

The model is the reason the question is worth asking. On LiveBench, a third-party benchmark that refreshes its
questions to limit contamination, the unquantized model at maximum reasoning effort ranks first of the
\dv{livebench.n_models} models listed on agentic coding (Section~\ref{sec:background}). Those scores belong to
DeepSeek's own service, not to the configuration we run, but they explain the interest. A model of that standing on
two workstation-class GPUs is a coding model that a company can rent for \dv{price.2xpro6000.eur_h}\,EUR an hour or
run on hardware it owns.

Table~\ref{tab:key-results} gives the answers a reader can take away. Each row points to the section that holds the
data and its caveats. Experiments carry the numbers E1 to E13 of the protocol in Section~\ref{sec:method}, and the
ones that did not run (among them E9, the reference with everything in GPU memory) are listed there. Every row
carries the same caveat. Each experiment ran on a single boot of the server on a single rented pod, so the intervals
describe repetitions within that boot and nothing measures how much a second boot or another pod would move the
numbers, and the cost of the offload against the model held entirely in GPU memory is still unmeasured.

\begin{table}[t]
\centering\small
\begin{tabular}{>{\raggedright\arraybackslash}p{3.3cm}>{\raggedright\arraybackslash}p{9.6cm}l}
\toprule
question & what we measured & section \\
\midrule
Does the model fit without the offload? & No. The pack needs \dv{static.noplan.vram_per_gpu_gib}\,GiB of weights per GPU
against \dv{expl.gpu_total_gib}\,GiB available. Both configurations without offload run out of memory while loading,
and the configuration with Diffbot's \texttt{gen30} plan boots and serves (E2). & \ref{sec:system} \\
What does a long prompt cost? & Reading it: a fresh 128K-token prompt takes
\dv{ctl.E12.long.prefill-long-131072.closed:1.ttft_p50}\,s to its first token. Writing barely changes for a single
request: \dv{ctl.E12.long.decode-256k.closed:1.tpot_p50_ms}\,ms per token after a 256K-token prompt against
\dv{ctl.E3.gen30-session2.A-code-4k.closed:1.tpot_p50_ms}\,ms after 4K tokens on the same pod (E12). &
\ref{sec:results-controlled} \\
How fast for one user, and for many? & One stream writes at a median of \dv{ctl.E8.stream.A.N1.p50} tokens/s. With a
cached 16K-token prefix, 64 users together get \dv{ctl.E8.main.B-prefix-16k.closed:64.agg_mean} tokens/s, but each
one only \dv{ctl.E8.stream.B.N64.p50} (E8). & \ref{sec:results-controlled} \\
Are the answers as good? & On fixed subsets, GSM8K gives the same outcome on every item as DeepSeek's API. The
differences on IFEval (\dv{ctl.Q.ifeval.diff}) and HumanEval (\dv{ctl.Q.humaneval_instruct.diff}) have 95\% intervals
that include zero (E10, E11). & \ref{sec:quality} \\
How many coding agents at once? & \dv{agent.E13.agent.Cprov} simultaneous agent-turn streams keep at least 40 output
tokens/s each with a 10\% margin, on a pod running nothing else (E13). & \ref{sec:agent} \\
How many developers per server? & About \dv{agent.E13.agent.seats_conc_e1} intermittently active developers, so that
fewer than one moment in a hundred has more than four of them waiting for the model, against
\dv{agent.E13.agent.seats_volume_u100} by monthly token volume at full utilization. A pod costs as much as
\dv{agent.break_even_seats} subscriptions (derived from E13). & \ref{sec:agent} \\
Can upstream vLLM do it alone? & Not on these pins. Its zero-copy offload of the same experts fails at initialization
with a KV block of 64 and on the first request with 128 (E3). & \ref{sec:results-controlled} \\
Does the host matter? & Yes. On a second host with a different NUMA layout and CPU, the median time per token was
\dv{ctl.pod.cross.A.N1.tpot_ratio} to \dv{ctl.pod.cross.A.N32.tpot_ratio} times higher (E3, one repetition per
pod). & \ref{sec:results-controlled} \\
\bottomrule
\end{tabular}
\caption{What this paper found. Pod measurements come from one boot per experiment on 2x RTX PRO 6000 pods, with the
model's step-by-step reasoning mode (thinking) turned off. The first row also rests on the byte accounting of the
checkpoint headers. The quality reference (E11) is DeepSeek's API on benchmark subsets, and the seat counts are
derived from E13 under the usage assumptions of Section~\ref{sec:agent}.}
\label{tab:key-results}
\end{table}

How many developers using a coding agent can share one server? The usual way to answer is by volume: count the output
tokens the server can produce in a month and divide by what one developer consumes. For the usage profile we measured,
that gives \dv{agent.E13.agent.seats_volume_u100} developers at full utilization and
\dv{agent.E13.agent.seats_volume_u50} at half. If each developer must see the model write at least 40 tokens per
second, the answer is \dv{agent.E13.agent.seats_conc_e1}, between \dv{agent.E13.agent.volume_over_conc_u50} and
\dv{agent.E13.agent.volume_over_conc_u100} times fewer. We designed the coding-agent experiment around this gap. The
server behaves like a bar that can pour for a whole office every night but has four stools at the counter, and
Section~\ref{sec:agent} explains why the number of stools decides how many regulars it can take.

\paragraph{How we measured.}
Before running anything we audited who wrote each part of the stack, from pinned sources and patches
(Section~\ref{sec:attribution}). We accounted every byte of both checkpoints from their safetensors headers, without
downloading the weights, to know what fits where before touching a GPU (Section~\ref{sec:system}). We then ran
experiments in three sessions on rented pods with the same GPU model: the boots without offload, a concurrency sweep,
two host shares and upstream vLLM's own offload, prompts up to 260,000 tokens, quality on three task subsets against
DeepSeek's API, and a coding-agent workload built from one developer's observed usage. Every number in the paper is
printed by a macro that reads a value computed from a data file, and the build fails if the file is missing
(Section~\ref{sec:repro}).

\paragraph{How to read this paper.}
Sections~\ref{sec:attribution} to~\ref{sec:discussion} give the findings in the order that makes them easiest to
follow: who built the stack, why the model does not fit and how it is made to, how fast the result runs, whether the
answers hold, what it means for coding agents, and what remains open. The record follows, with the earlier
exploratory runs (Section~\ref{sec:results-exploratory}), the full methodology, the setup of the pods, related work,
limitations and reproducibility. A reader who wants to check a claim can go from any sentence of the first part to
the protocol and the data file behind it.

\paragraph{What this paper does not claim.}
It claims no new algorithm or system. The offload is third-party work, and our part is the integration, the
configuration and the measurement. It does not claim that the stack preserves the model's quality in general, since
our quality evidence covers three small task subsets with thinking off. And it does not yet say what the offload costs
against the same model held entirely in GPU memory, because that reference needs larger hardware and was not run.
